% English translation of chapter7.tex lines 892--1129.
% Source: Methods of Algebra, Volume 2, commit
% 9a5803ff77dd3257484cb177f851a73770a59dd3 (CC BY 4.0).
% stable-unit-id: o014.aljabr2.chapter7.coalgebra-to-hopf-a
% Coverage: Section 7.5 from its opening through the proof of the monoidal
% structure on the module category of a bialgebra; continued in Unit 093.

% segment-id: o014.li2.chapter7.coalgebra-to-hopf-a.h001
\section{From Coalgebras to Hopf Algebras}\label{sec:bialgebra}
\hypertarget{o014.aljabr2.chapter7.coalgebra-to-hopf-a}{ }
% segment-id: o014.li2.chapter7.coalgebra-to-hopf-a.p001
The first aim of this section is to dualize
Definition~\ref{def:algebra-monoidal} of an algebra. To this end, first note
that the definition of a monoidal category is self-dual. Let $\mathcal{V}$ be
a monoidal category. On $\mathcal{V}^{\opp}$, one can retain the original
bifunctor $\otimes$ and unit object $\munit$, but replace the associativity
constraint $a$ and $\iota:\munit\otimes\munit\rightiso\munit$ in
\cite[Definition 3.1.1]{Li1} by their inverses. In this way,
$\mathcal{V}^{\opp}$ becomes a monoidal category.

% segment-id: o014.li2.chapter7.coalgebra-to-hopf-a.p002
Similarly, if $\mathcal{V}$ has a braiding $c$ as in
\cite[Definition 3.3.1]{Li1}, taking inverses gives the corresponding
braiding on $\mathcal{V}^{\opp}$. Under this correspondence, $\mathcal{V}$ is
a symmetric monoidal category if and only if $\mathcal{V}^{\opp}$ is one.

% segment-id: o014.li2.chapter7.coalgebra-to-hopf-a.d001
\begin{definition}[Coalgebra]\label{def:cogebra-monoidal}
	\index{coalgebra}
	Let $\mathcal{V}$ be a monoidal category. An algebra over
	$\mathcal{V}^{\opp}$ is called a \emph{coalgebra} over $\mathcal{V}$. In
	other words, a coalgebra is data $(C,\Delta,\epsilon)$, where
	% segment-id: o014.li2.chapter7.coalgebra-to-hopf-a.q001
	\[ C \in \Obj(\mathcal{V}), \quad \Delta = \Delta_C: C \to C \otimes C, \quad \epsilon = \epsilon_C: C \to \munit, \]
	satisfying the commutative diagrams dual to those in
	Definition~\ref{def:algebra-monoidal}:
	% segment-id: o014.li2.chapter7.coalgebra-to-hopf-a.q002
	\[\begin{tikzcd}
			\munit \otimes C & C \otimes C \arrow[r, "{\identity \otimes \epsilon}"] \arrow[l, "{\epsilon \otimes \identity}"'] & C \otimes \munit \\
			& C \arrow[lu, "\sim"' sloped] \arrow[ru, "\sim"' sloped] \arrow[u, "\Delta"'] &
		\end{tikzcd} \quad
	\begin{tikzcd}[column sep=tiny]
		(C \otimes C) \otimes C & & C \otimes (C \otimes C) \arrow[ll, "\sim"] \\
		C \otimes C \arrow[u, "{\Delta \otimes \identity}"] & C \arrow[l, "\Delta"] \arrow[r, "\Delta"'] & C \otimes C . \arrow[u, "{\identity \otimes \Delta}"']
	\end{tikzcd}\]

	The map $\Delta$ is customarily called the \emph{comultiplication} on $C$,
	and $\epsilon$ the \emph{counit} of $C$.

	A morphism from a coalgebra $(C,\Delta,\epsilon)$ to
	$(C',\Delta',\epsilon')$ is defined dually to
	Definition~\ref{def:algebra-monoidal}: it is a morphism $\phi:C\to C'$ for
	which the following diagrams commute:
	% segment-id: o014.li2.chapter7.coalgebra-to-hopf-a.q003
	\[\begin{tikzcd}
		\munit & C \arrow[d, "\phi"] \arrow[l, "\epsilon"']  \\
		& C' \arrow[lu, "{\epsilon'}"]
	\end{tikzcd} \quad \begin{tikzcd}
		C \otimes C \arrow[r, "{\phi \otimes \phi}"] & C' \otimes C' \\
		C \arrow[r, "\phi"'] \arrow[u, "\Delta"] & C' . \arrow[u, "{\Delta'}"']
	\end{tikzcd}\]

	If $\mathcal{V}$ is equipped with a braiding $c$ and $c\Delta=\Delta$,
	then $(C,\Delta,\epsilon)$ is called \emph{cocommutative}.
\end{definition}

% segment-id: o014.li2.chapter7.coalgebra-to-hopf-a.p003
The following definition is simply the dual of
Definition~\ref{def:module-algebra}.

% segment-id: o014.li2.chapter7.coalgebra-to-hopf-a.d002
\begin{definition}[Comodule]\label{def:comodule-cogebra}
	\index{comodule}
	Let $(C,\Delta,\epsilon)$ be a coalgebra over $\mathcal{V}$. A left
	$C$-comodule is defined to be data $(M,\rho)$, where
	% segment-id: o014.li2.chapter7.coalgebra-to-hopf-a.q004
	\[ M \in \Obj(\mathcal{V}), \quad \rho = \rho_M : M \to C \otimes M, \]
	and $\rho$ may be viewed as scalar comultiplication on the left comodule.
	The condition is that the following diagrams commute:
	% segment-id: o014.li2.chapter7.coalgebra-to-hopf-a.q005
	\[\begin{tikzcd}
		\munit \otimes M & C \otimes M \arrow[l, "{\epsilon \otimes \identity}"'] \\
		& M \arrow[u, "{\rho}"'] \arrow[lu, "\sim" sloped]
	\end{tikzcd} \quad \begin{tikzcd}[column sep=small]
		(C \otimes C) \otimes M & & C \otimes (C \otimes M) \arrow[ll, "\sim"'] \\
		C \otimes M \arrow[u, "{\Delta \otimes \identity}"] & M \arrow[l, "{\rho}"] \arrow[r, "{\rho}"'] & C \otimes M. \arrow[u, "{\identity \otimes \rho}"']
	\end{tikzcd}\]
	A homomorphism from a left comodule $(M,\rho)$ to $(M',\rho')$ is a
	morphism $f:M\to M'$ satisfying
	$(\identity\otimes f)\rho=\rho'f$. Right $C$-comodules, bicomodules, and
	homomorphisms between them are defined similarly.
\end{definition}

% segment-id: o014.li2.chapter7.coalgebra-to-hopf-a.p004
An elementary but important example is $C$ itself: with
$\rho_C:=\Delta:C\to C\otimes C$, it becomes both a left and a right
comodule.

% segment-id: o014.li2.chapter7.coalgebra-to-hopf-a.p005
Although the definition of a comodule may seem contrary to one's usual
intuition, working with comodules need not be more complicated than working
with modules. For example, take $\mathcal{V}=\Bbbk\dcate{Mod}$, where
$\Bbbk$ is a commutative ring. Let $M$ be a free $\Bbbk$-module with basis
$(v_i)_{i\in I}$. Specifying $\rho:M\to M\dotimes{\Bbbk}C$ is equivalent to
specifying a family of elements $(t_{ij})_{(i,j)\in I^2}$ of $C$ such that,
for each $i\in I$,
% segment-id: o014.li2.chapter7.coalgebra-to-hopf-a.q006
\begin{equation*}
	\rho(v_i) = \sum_{j \in I} v_j \otimes t_{ji} \quad \text{(finite sum)},
\end{equation*}
and the comodule conditions are expressed, for every $i$, by
% segment-id: o014.li2.chapter7.coalgebra-to-hopf-a.q007
\begin{equation*}
	v_i = \sum_j \epsilon(t_{ji}) v_j, \quad \sum_j v_j \otimes \Delta(t_{ji}) = \sum_{j, k} v_k \otimes t_{kj}  \otimes t_{ji}.
\end{equation*}
Renaming the indices appropriately, these conditions simplify further, for
all $i,j$, to
% segment-id: o014.li2.chapter7.coalgebra-to-hopf-a.q008
\begin{equation}\label{eqn:comodule-matrix}
	\begin{aligned}
		\epsilon(t_{ji}) & = \delta_{j, i}, \\
		\Delta(t_{ji}) & = \sum_k t_{jk} \otimes t_{ki};
	\end{aligned}
\end{equation}
here $\delta_{j,i}$ is the Kronecker delta, defined to be $1$ when $j=i$ and
$0$ otherwise.
\index{Kronecker delta}
\index[sym1]{delta-ij@$\delta_{ij}$}

% segment-id: o014.li2.chapter7.coalgebra-to-hopf-a.p006
Continuing the preceding discussion, express a $\Bbbk$-linear map
$f:M\to M'$ between right comodules relative to the chosen bases as
$f(v_i)=\sum_jr_{ji}v'_j$, that is, as a matrix with $r_{ji}\in\Bbbk$. A
similar argument proves that $f$ is a comodule homomorphism if and only if,
for all $i,k$,
% segment-id: o014.li2.chapter7.coalgebra-to-hopf-a.q009
\begin{equation*}\begin{gathered}
	\sum_j r_{kj} t_{ji} = \sum_j r_{ji} t'_{kj},
\end{gathered}\end{equation*}
where $\Bbbk$ acts on $C$ from the left by scalar multiplication. The version
for left comodules is, of course, similar.

% segment-id: o014.li2.chapter7.coalgebra-to-hopf-a.p007
We shall continue the discussion of several aspects of comodule theory in
\S\ref{sec:Morita}.

% segment-id: o014.li2.chapter7.coalgebra-to-hopf-a.r001
\begin{remark}\label{rem:coalgebra-convolution}
	\index{convolution}
	If $(C,\Delta,\epsilon)$ is a coalgebra over $\mathcal{V}$ and
	$(A,\mu,\eta)$ is an algebra over $\mathcal{V}$, then $\Hom(C,A)$ is
	equipped with a binary operation $\star$, called \emph{convolution}, as
	follows:
	% segment-id: o014.li2.chapter7.coalgebra-to-hopf-a.q010
	\[ f \star g := \mu (f \otimes g) \Delta, \quad f, g \in \Hom(C, A). \]
	The properties of $\Delta$ and $\mu$ readily show that $\star$ is
	associative, and the following commutative diagram suffices to show that
	$(\Hom(C,A),\star)$ is a monoid with unit
	$\eta\epsilon\in\Hom(C,A)$:
	% segment-id: o014.li2.chapter7.coalgebra-to-hopf-a.q011
	\[\begin{tikzcd}[column sep=large]
		A \otimes \munit \arrow[r, "{\mu(\identity \otimes \eta)}"] & A & \munit \otimes A \arrow[l, "{\mu(\eta \otimes \identity)}"'] \\
		C \otimes \munit \arrow[r] \arrow[u, "{f \otimes \identity}"] & C \arrow[u, "f"] & \munit \otimes C \arrow[l] \arrow[u, "{\identity \otimes f}"'] \\
		C \otimes C \arrow[u, "{\identity \otimes \epsilon}"] & C \arrow[l, "\Delta"] \arrow[r, "\Delta"'] \arrow[equal, u] & C \otimes C . \arrow[u, "{\epsilon \otimes \identity}"']
	\end{tikzcd}\]
	A standard argument shows that if $\varphi:C'\to C$ is a coalgebra
	morphism and $\psi:A\to A'$ an algebra morphism, then
	% segment-id: o014.li2.chapter7.coalgebra-to-hopf-a.q012
	\begin{equation}\label{eqn:convolution-functorial}
		\Hom(C, A) \to \Hom(C', A'), \quad f \mapsto \psi f \varphi
	\end{equation}
	is a homomorphism between the convolution monoids.
\end{remark}

% segment-id: o014.li2.chapter7.coalgebra-to-hopf-a.p008
From now on, let $\mathcal{V}$ be a braided monoidal category, and write its
braiding in the usual form
% segment-id: o014.li2.chapter7.coalgebra-to-hopf-a.q013
\[ c(X, Y): X \otimes Y \to Y \otimes X. \]
By \S\ref{sec:alg-in-monoidal-cat}, the category of algebras
$\cate{Alg}(\mathcal{V})$ (or the category of coalgebras
$\cate{Alg}(\mathcal{V}^{\opp})^{\opp}$) has a monoidal structure with unit
$\munit$. We may therefore consider coalgebras in the category of algebras
(or algebras in the category of coalgebras). A simple but interesting
observation is that both lead to the same data
$(A,\mu,\eta,\Delta,\epsilon)$, with $A\in\Obj(\mathcal{V})$. The conditions
required of the morphisms
$A\xrightarrow{\Delta}A\otimes A\xrightarrow{\mu}A$ and
$\munit\xrightarrow{\eta}A\xrightarrow{\epsilon}\munit$ in $\mathcal{V}$
are:
% segment-id: o014.li2.chapter7.coalgebra-to-hopf-a.l001
\begin{enumerate}[(i)]
	\item $(A,\mu,\eta)$ is an algebra;
	\item $(A,\Delta,\epsilon)$ is a coalgebra;
	\item the comultiplication $\Delta:A\to A\otimes A$ and counit
	$\epsilon:A\to\munit$ are algebra morphisms;
	\item the multiplication $\mu:A\otimes A\to A$ and unit
	$\eta:\munit\to A$ are coalgebra morphisms.
\end{enumerate}
Once (i) and (ii) are assumed, (iii) and (iv) are equivalent, so it is enough
to check either one. Both state the compatibility of $(\Delta,\epsilon)$ and
$(\mu,\eta)$ through four commutative diagrams, one for each possible pair;
in formulas,
% segment-id: o014.li2.chapter7.coalgebra-to-hopf-a.q014
\begin{equation*}\begin{gathered}
	\Delta \mu = (\mu \otimes \mu)(\identity_A \otimes c(A, A) \otimes \identity_A) (\Delta \otimes \Delta), \\
	\Delta \eta = \eta \otimes \eta, \quad \epsilon \mu = \epsilon \otimes \epsilon, \quad \epsilon \eta = \identity_{\munit}.
\end{gathered}\end{equation*}

% segment-id: o014.li2.chapter7.coalgebra-to-hopf-a.d003
\begin{definition}[Bialgebra]
	\index{bialgebra}
	Let $\mathcal{V}$ be a braided monoidal category. If $A$ is an algebra in
	the category of coalgebras over $\mathcal{V}$, or equivalently a coalgebra
	in the category of algebras over $\mathcal{V}$, then $A$ is called a
	\emph{bialgebra} over $\mathcal{V}$. Concretely, a bialgebra consists of
	the data $(A,\mu,\eta,\Delta,\epsilon)$.

	If $A$ and $A'$ are both bialgebras and $\phi:A\to A'$ is simultaneously
	an algebra morphism and a coalgebra morphism, then $\phi$ is called a
	morphism from the bialgebra $A$ to the bialgebra $A'$.
\end{definition}

% segment-id: o014.li2.chapter7.coalgebra-to-hopf-a.p009
The full data of a bialgebra are often abbreviated simply as $A$.

% segment-id: o014.li2.chapter7.coalgebra-to-hopf-a.p010
If $\Bbbk$ is a commutative ring and
$\mathcal{V}=\Bbbk\dcate{Mod}$, then algebras, coalgebras, and bialgebras
over $\mathcal{V}$ are also called $\Bbbk$-algebras, $\Bbbk$-coalgebras,
and $\Bbbk$-bialgebras, respectively, and so forth.

% segment-id: o014.li2.chapter7.coalgebra-to-hopf-a.eg001
\begin{example}\label{eg:monoid-bialgebra}
	Let $M$ be a monoid and $\Bbbk$ a commutative ring. Form the monoid algebra
	$\Bbbk[M]$ over $\Bbbk$; see \cite[Definition 5.6.1]{Li1}. Define the
	$\Bbbk$-module homomorphism
	% segment-id: o014.li2.chapter7.coalgebra-to-hopf-a.q015
	\begin{align*}
		\Delta: \Bbbk[M] & \to \Bbbk[M] \dotimes{\Bbbk} \Bbbk[M] \\
		m & \mapsto m \otimes m, \quad m \in M
	\end{align*}
	and $\epsilon:\Bbbk[M]\to\Bbbk$ by mapping every $m\in M$ to $1$. Let
	us verify that these data give a $\Bbbk$-bialgebra.

	First, $(\Bbbk[M],\Delta,\epsilon)$ is a coalgebra, since for every
	$m\in M$,
	% segment-id: o014.li2.chapter7.coalgebra-to-hopf-a.q016
	\begin{gather*}
		(\identity \otimes \epsilon)(\Delta(m)) = m = (\epsilon \otimes \identity)(\Delta(m)), \\
		(\Delta \otimes \identity)(\Delta(m)) = m \otimes m \otimes m = (\identity \otimes \Delta)(\Delta(m)).
	\end{gather*}
	Note that $\Bbbk[M]$ is always cocommutative, whereas $\Bbbk[M]$ is
	commutative if and only if $M$ is commutative.

	Second, check that $\epsilon$ and $\Delta$ are both homomorphisms of
	$\Bbbk$-algebras. It is enough to verify, for every $m,m'\in M$, that
	% segment-id: o014.li2.chapter7.coalgebra-to-hopf-a.q017
	\begin{gather*}
		\epsilon(mm') = 1 = \epsilon(m)\epsilon(m'), \quad \epsilon(1) = 1, \\
		\Delta(mm') = mm' \otimes mm' = (m \otimes m)(m' \otimes m') = \Delta(m) \Delta(m'), \\
		\Delta(1) = 1 \otimes 1.
	\end{gather*}
\end{example}

% segment-id: o014.li2.chapter7.coalgebra-to-hopf-a.p011
For an algebra $A$ over $\mathcal{V}$,
Definition~\ref{def:module-algebra} describes left $A$-modules. If $A$ is a
bialgebra, we can go further as follows.
% segment-id: o014.li2.chapter7.coalgebra-to-hopf-a.l002
\begin{enumerate}
	\item On the unit object $\munit$, define a left $A$-module structure by
	taking $\mu_{\munit}:A\otimes\munit\to\munit$ to be the composite
	$A\otimes\munit\rightiso A\xrightarrow{\epsilon}\munit$.
	\item For left $A$-modules $M_i$, denote the corresponding scalar
	multiplication morphisms by $\mu_i$ ($i=1,2$). Define the morphism
	$\mu_{M_1\otimes M_2}:A\otimes(M_1\otimes M_2)\to M_1\otimes M_2$ to be
	the composite
	% segment-id: o014.li2.chapter7.coalgebra-to-hopf-a.q018
	\begin{multline*}
		A \otimes M_1 \otimes M_2 \xrightarrow{\Delta \otimes \identity_{M_1 \otimes M_2}} A \otimes A \otimes M_1 \otimes M_2 \\
		\xrightarrow{\identity_A \otimes c(A, M_1) \otimes \identity_{M_2}} A \otimes M_1 \otimes A \otimes M_2 \xrightarrow{\mu_1 \otimes \mu_2} M_1 \otimes M_2.
	\end{multline*}
	Here, to simplify notation, parentheses in $\otimes$ operations and the
	associativity constraints have been suppressed; in other words,
	$\mathcal{V}$ is treated as a strict monoidal category.
\end{enumerate}

% segment-id: o014.li2.chapter7.coalgebra-to-hopf-a.p012
These definitions are intended to give the objects $\munit$ and
$M_1\otimes M_2$ of $\mathcal{V}$ left $A$-module structures. The case of
right $A$-modules is, of course, similar.

% segment-id: o014.li2.chapter7.coalgebra-to-hopf-a.pr001
\begin{proposition}\label{prop:bialgebra-Mod-monoidal}
	\index[sym1]{Comod@$\cate{Comod}$}
	Let $A$ be a bialgebra over a braided monoidal category $\mathcal{V}$. The
	definitions above make the category of left $A$-modules $A\dcate{Mod}$ a
	monoidal category with unit $(\munit,\mu_{\munit})$. The same holds for the
	category of right $A$-modules $\cated{Mod}A$.

	Dually, the multiplication and unit of $A$ make the category of left
	$A$-comodules $A\dcate{Comod}$ and the category of right $A$-comodules
	$\cated{Comod}A$ into monoidal categories.

	For all four categories above, the forgetful functor to $\mathcal{V}$ is
	always a monoidal functor.
\end{proposition}
% segment-id: o014.li2.chapter7.coalgebra-to-hopf-a.pf001
\begin{proof}
	By duality, it is enough to prove the first assertion. The verification is
	lengthy but straightforward; we outline it. First, one must check that
	scalar multiplication satisfies the commutative diagrams in
	Definition~\ref{def:module-algebra}. For $(\munit,\mu_{\munit})$, these
	properties follow from the fact that $\epsilon:A\to\munit$ is an algebra
	morphism. For the module structure on $M_1\otimes M_2$, the diagram for
	scalar multiplication by the unit follows from the fact that
	$\Delta:A\to A\otimes A$ is an algebra morphism, in particular from its
	compatibility with $\eta$ and $\eta\otimes\eta$.

	For associativity of scalar multiplication, after expanding the
	definition, the issue is to prove that two maps from
	$A\otimes A\otimes M_1\otimes M_2$ to $M_1\otimes M_2$ are equal:
	% segment-id: o014.li2.chapter7.coalgebra-to-hopf-a.l003
	\begin{enumerate}[(a)]
		\item first apply
		$\Delta\otimes\Delta\otimes\identity_{M_1}\otimes\identity_{M_2}$,
		then, taking the braiding into account, successively let the first and
		third (or second and fourth) copies of $A$ in
		$A\otimes A\otimes(A\otimes A\otimes M_1\otimes M_2)$ act from the
		left on $M_1$ (or $M_2$);
		\item first apply $\mu:A\otimes A\to A$ in the first two positions,
		then apply $\mu_{M_1\otimes M_2}$. Since
		% segment-id: o014.li2.chapter7.coalgebra-to-hopf-a.q019
		\[ \Delta\mu = (\mu \otimes \mu) (\identity \otimes c(A, A) \otimes \identity)(\Delta \otimes \Delta), \]
		this map is the composite
		% segment-id: o014.li2.chapter7.coalgebra-to-hopf-a.q020
		\begin{multline*}
			A \otimes A \otimes M_1 \otimes M_2 \xrightarrow{\Delta \otimes \Delta \otimes \identity_{M_1} \otimes \identity_{M_2}} A \otimes A \otimes A \otimes A \otimes M_1 \otimes M_2 \\
			\xrightarrow{\identity_A \otimes c(A, A) \otimes \identity_A \otimes \identity_{M_1} \otimes \identity_{M_2}} \boxed{A \otimes A \otimes A \otimes A} \otimes M_1 \otimes M_2 \\
			\xrightarrow{\mu \otimes \mu \otimes \identity_{M_1} \otimes \identity_{M_2}} A \otimes A \otimes M_1 \otimes M_2 \\
			\xrightarrow{\identity_A \otimes c(A, M_1) \otimes \identity_{M_2}} A \otimes M_1 \otimes A \otimes M_2 \xrightarrow{\mu_1 \otimes \mu_2} M_1 \otimes M_2.
		\end{multline*}
		The scalar multiplications $\mu_1$ and $\mu_2$ are each associative and
		$c$ is functorial. Thus, taking the braiding into account, the last
		three arrows in this composite may also be regarded as first letting
		the first and second (or third and fourth) copies of $A$ in the boxed
		part act successively from the left on $M_1$ (or $M_2$).
	\end{enumerate}

	Braid diagrams help explain this, as at the end of
	\cite[\S 7.4]{Li1}. To prove that (a) equals (b), it is enough to verify
	the following equality\footnote{There is one harmless minor difference:
	in the cited reference, morphisms are composed from bottom to top, whereas
	here they are composed in the opposite direction.}
	% segment-id: o014.li2.chapter7.coalgebra-to-hopf-a.g001
	\begin{center}\begin{tikzpicture}
		[baseline=(braid), yscale=0.8]
		\pic[
			braid/number of strands = 6,
			braid/every strand/.style={black},
			name prefix=braid
		] at (0, 0) {braid={s_4 s_2 s_3}};
		\coordinate (braid) at ($(braid-1-s)!0.5!(braid-1-e)$);
		\node[above] at (braid-1-s) {$A$};
		\node[below=0.3em] (P) at (braid-1-e) {$A$};
		\node[above] at (braid-2-s) {$A$};
		\node[below=0.3em] (R) at (braid-2-e) {$A$};
		\node[above] at (braid-3-s) {$A$};
		\node[below=0.3em] at (braid-3-e) {$A$};
		\node[above] at (braid-4-s) {$A$};
		\node[below=0.3em] at (braid-4-e) {$A$};
		\node[above] at (braid-5-s) {$M_1$};
		\node[below=0.3em] (Q) at (braid-5-e) {$M_1$};
		\node[above] at (braid-6-s) {$M_2$};
		\node[below=0.3em] (S) at (braid-6-e) {$M_2$};

		\node (X) [draw, fit=(P) (Q)] {};
		\node[below=0.2em of X] {act on\; $M_1$};
		\node (Y) [draw, fit=(R) (S)] {};
		\node[below=0.2em of Y] {act on\; $M_2$};
	\end{tikzpicture} \; = \begin{tikzpicture}
		[baseline=(braid), yscale=0.8]
		\pic[
			braid/number of strands=6,
			braid/every strand/.style = {black},
			name prefix=braid
		] at (0, 0) {braid={s_2 s_4 s_3}};
		\coordinate (braid) at ($(braid-1-s)!0.5!(braid-1-e)$);
		\node[above] at (braid-1-s) {$A$};
		\node[below=0.3em] (P) at (braid-1-e) {$A$};
		\node[above] at (braid-2-s) {$A$};
		\node[below=0.3em] (R) at (braid-2-e) {$A$};
		\node[above] at (braid-3-s) {$A$};
		\node[below=0.3em] at (braid-3-e) {$A$};
		\node[above] at (braid-4-s) {$A$};
		\node[below=0.3em] at (braid-4-e) {$A$};
		\node[above] at (braid-5-s) {$M_1$};
		\node[below=0.3em] (Q) at (braid-5-e) {$M_1$};
		\node[above] at (braid-6-s) {$M_2$};
		\node[below=0.3em] (S) at (braid-6-e) {$M_2$};

		\node (X) [draw, fit=(P) (Q)] {};
		\node[below=0.2em of X] {act on\; $M_1$};
		\node (Y) [draw, fit=(R) (S)] {};
		\node[below=0.2em of Y] {act on\; $M_2$};
	\end{tikzpicture}\end{center}
	which is immediately apparent.

	To show that $A\dcate{Mod}$ is a monoidal category, one must still prove
	that the associativity and unit constraints for $\otimes$ are all
	morphisms of left $A$-modules. The first assertion is straightforward. The
	second requires the coalgebra property
	$(\identity\otimes\epsilon)\Delta=\identity_A=
	(\epsilon\otimes\identity)\Delta$, after identifying
	$A\otimes\munit\simeq A\simeq\munit\otimes A$.
\end{proof}
